\section{Lessons}

Throughout this project, we have learned and practiced various essential steps in the Machine Learning workflow, which helped deepen our understanding of both theoretical concepts and hands-on implementations. The key lessons learned are summarized as follows:

\begin{itemize}
    \item \textbf{Data Collection and Gathering}: 
    We collected and gathered the necessary datasets, ensuring that the data sources were reliable and the data was suitable for the problem we aimed to solve.

    \item \textbf{Data Preprocessing and Normalization}: 
    We performed preprocessing steps including handling missing values, encoding categorical data, and applying normalization techniques to scale features properly. This step helped improve the stability and performance of our models.

    \item \textbf{Data Visualization}: 
    We utilized various visualization techniques to explore the data distributions and relationships between features. This helped in gaining insights into the dataset and detecting potential issues such as imbalance or outliers.

    \item \textbf{Building Features}: 
    Feature engineering was applied to enhance the model's learning capability. We built features using different vectorization techniques such as:
    \begin{itemize}
        \item \textbf{Count Vectorization}
        \item \textbf{TF-IDF (Term Frequency-Inverse Document Frequency)}
        \item \textbf{Word2Vec}
        \item \textbf{GloVe (Global Vectors for Word Representation)}
    \end{itemize}
    These approaches allowed us to transform raw text data into meaningful numerical representations.

    \item \textbf{Model Selection}: 
    We explored and implemented a variety of models, including:
    \begin{itemize}
        \item \textbf{Graphical Models}: Bayesian Networks, Hidden Markov Models (HMM)
        \item \textbf{Support Vector Machines (SVM)}: with different kernels and regularization techniques
        \item \textbf{PCA for Dimension Reduction}: with feature selection strategies, topic modeling, variance analysis, and chi-squared testing
        \item \textbf{LDA as Linear Classifier}
        \item \textbf{Ensemble Methods}: Random Forest, XGBoost, Stacking, and Voting strategies using Logistic Regressions and other selected models
        \item \textbf{Discriminative Models}: Logistic Regression, BiLSTM-CRF (to leverage the superior performance of the BiLSTM architecture)
    \end{itemize}

    \item \textbf{Training, Evaluating, and Testing}: 
    We trained the models using appropriate training-validation splits and evaluated them using consistent metrics. Hyperparameter tuning and model validation helped ensure the robustness of our solutions.

    \item \textbf{Performance Comparison}: 
    We performed a comprehensive comparison between the models, analyzing their strengths, weaknesses, and suitability for the given task. This analysis guided us in selecting the best-performing approach.
\end{itemize}
